%=====================================================================
\section{Analysis}
\label{sec:analysis}

\paragraph{Options must follow the annotation, not the inventory.}
The revised \umr procedure moved every question into data files and completed its inventories from the guidelines, among them the 27 aspect values of the UMR aspect lattice. The gold annotation uses seven of them. Given the full lattice, Jev chose fine-grained values the gold never uses (\emph{directed-achievement} 118 times, \emph{irreversible-directed-achievement} 59 times on the Chinese test set), the aspect score of the scorer's fine-grained judge fell from 0.70 for the procedure it revised to 0.31, and F fell to 33.3. Restricting the options to the seven values of the gold (aspect 0.71), together with one polarity rule the revision had dropped, gave 38.5 (Table~\ref{tab:umr}). A decision model answers the question it is asked; whether the options are the right ones is a property of the procedure, and only measurement against gold reveals it. In the same spirit, we tried asking aspect in two steps (a coarse kind, then the value within it): with the same decision model and test set, aspect accuracy dropped from 0.72 to 0.67 and F by 0.3, so we kept the single question.

\paragraph{Where UMR loses: structure, not labels alone.}
To see what closed decisions leave on the table in \umr, we replaced parts of K4's test output by gold, using the scorer's alignment of predicted to gold nodes, and re-scored. Making every concept form, aspect, role label and attribute of the aligned nodes correct raises F from 40.4 to 49.0; additionally making the core and peripheral arguments among aligned nodes correct raises it to 55.9. The remaining gap is nodes that the procedure never built. Correct aspect values alone are worth 2.9 F, but they are the least learnable part: a lexical prior from the training gold does not beat Jev on the test set (0.714 vs.\ at most 0.716 accuracy), since the same verbs receive different aspects across documents. Attachment is the procedure's own work: rules split a sentence top-down and assign each word to a role by a per-word question, and an early routing error moves a whole subtree. On PARSEME, by contrast, the structure is a set of token groups proposed by rules, and every remaining error is a decision or a missing candidate rule, which is exactly what reflection edits.

\paragraph{The development set decides whether reflection transfers.}
The three kinds of evidence a reflection round uses (the batch, earlier batches and the development set) protect against different failures. On PARSEME, with 150 development sentences, development gains carried over to the test set (+7.6 dev, +5.5 test in Chinese). On Chinese \umr, a first run (from the earlier seed J0) with a ten-sentence development set accepted 18 of 32 edits and raised development F by 7.3 but test F by only 1.6. A second run (from $\Pi_0$) with 27 development sentences accepted 15 of 42 edits; development F rose by 0.6 and test F did not change. The seeds differ, so the comparison is suggestive, but the pattern is the expected one. A small development set accepts edits that fit its own sentences; a larger one rejects them, which is the intended behavior but leaves the batches' gains behind. Measurement makes this visible: every kept edit has a recorded effect on each set.

\paragraph{What the procedure costs.}
Decisions are the cheap part. Deciding the 3,462 sentences of the Chinese PARSEME test set cost \$0.06; a \umr sentence takes 11--49 requests to Jev and costs \$0.001--0.006, depending on the procedure. Reflection costs a few dollars per run and is paid once. Routing costs more per decision than Jev by two orders of magnitude, but only for uncertain questions (Chinese PARSEME: \$0.09 for the test set). Direct generation by a cheap LLM costs \$0.11--1.16 per test set, more than the procedure's decisions on PARSEME and less than the procedure on \umr.
